\section{Научная новизна}

Предложенный подход вносит следующий вклад в область обнаружения кибератак:

\begin{enumerate}
    \item \textbf{Формализация атаки как иерархического атрибутивного метаграфа}, объединяющего стратегические цели, тактики (MITRE ATT\&CK) и техники реализации в единую семантическую структуру с гиперсвязями. Это принципиально отличает модель от деревьев атак и классических графов. 
    \item \textbf{Алгоритм частичного изоморфизма с онтологически обоснованной совместимостью атрибутов}, устойчивый к обфускации и LotL-техникам.    
    \item \textbf{Интеграция регуляторных требований (ФСТЭК, ГОСТ Р 56939–2024)} непосредственно в модель через атрибуты, что обеспечивает соответствие отечественным стандартам защиты информации.
\end{enumerate}

В работе впервые предложена \textbf{онтологически насыщенная, регуляторно-ориентированная модель выявления атак}, основанная на структурно-семантическом анализе атрибутивных метаграфов.